\subsection{KV Master: Workflow-aware Tiered Placement}
\label{sec:kv-master}

Master first combines request urgency with within-request KV demand,
then carries this priority across repeated and shared uses. Renewable
credit manages lower-tier retention; a fixed-size, request-level policy
allocates Peer acceleration. An event loop connects both policies to
actual transfers and available capacity.

\subsubsection{Slack- and Access-aware KV Priority}
\label{sec:placement-value}
\label{sec:prefix-value}

Two factors determine a request's interest in a KV block:
how strongly recovery affects its workflow, and how likely execution is
to need that block. Let $R_r$ be the profiled recovery delay under the
request's current placement, including lower-tier reads, replay and
queues, and $s_r$ its guarded slack. We use the exposed fraction
\begin{equation}
 u_r=\frac{[R_r-s_r]^+}{\max(R_r,\varepsilon)},
 \qquad P_{r,l,b}=u_r\,\widehat p_{r,l,b},
 \label{eq:value}
\end{equation}
where $\varepsilon>0$ avoids division by zero. Thus zero slack gives full
urgency when recovery takes time; slack covering recovery removes this
urgency, not ordinary reuse value. Invalid Tracker estimates use zero
slack. $P$ is a ranking signal, not an additive estimate of WCT saved.

Within one dense-attention layer, every causally visible historical block
is needed by subsequent queries, so its access factor is one. Earlier
creation alone does not imply greater future value. Sparse attention
instead uses a recent window of completed queries:
\begin{equation}
 \widehat p_{r,l,b}=
 \frac{\text{queries selecting block }b}
      {\text{queries for which block }b\text{ was visible}}.
 \label{eq:selection-probability}
\end{equation}
Each query counts a physical block once across heads. Visibility is
defined before sparse selection, not after candidate pruning. A block
selected 40 times in 200 opportunities has lower estimated demand than
one selected 15 times in 30; raw counts reverse that ranking. A calibrated
layer prior smooths small denominators. This causal estimate predicts
demand without modifying the native selector; mandatory auxiliary state
remains a dependency regardless of its measured frequency.

Priority and retention cost are distinct. For equal-size, equal-cost
blocks in the same dense layer, the access contribution is equal; object
sizes, transfer paths and reconstruction costs can still differ.
Figure~\ref{fig:retention-value} illustrates why cost per byte also matters.
This motivates cost-aware lower-tier credit, while the Peer policy below
deliberately uses a simpler urgency rule.

\begin{figure}[t]
 \centering
 \includegraphics[width=\columnwidth]{figures/retention-value.pdf}
 \caption{Slack alone does not rank retention value. In this isolated-fault
 example, A uses two capacity units to reduce recovery from 4 to 2\,s;
 B uses one to reduce 18 to 8\,s. A at slack 0 saves 1\,s/unit, whereas
 B at slack 6 saves 10\,s/unit. Normal-use value is zero. This constructed
 example motivates cost-aware lower-tier retention and exposes a limit
 of urgency-only Peer admission; it is not a policy-speedup result.}
 \label{fig:retention-value}
 \Description{A higher-slack object can save more workflow time per byte.}
\end{figure}

\subsubsection{Renewable Credit with Lazy Updates}
\label{sec:renewable-credit}

The preceding priority concerns one request. Across requests, a shared
prefix can serve several Agents, while one Agent can return after a tool
wait. Master gives each interested Agent its own share of an object's
credit. A different Agent adds a share; the same Agent's new use refreshes
its share instead of accumulating another copy of its value. Agent
identity includes its workflow instance, and physical capacity is charged
only once.

For a lower-tier object $x$, the current grant combines a normal-reuse
term with its recovery priority:
\begin{equation}
 g_{a,x,j}=v^{\mathrm{reuse}}_{a,x,j}
                +\lambda_f c_{a,x,j}P_{a,x}.
 \label{eq:credit-grant}
\end{equation}
Here $j$ is DRAM or storage; $\lambda_f$ weights fault risk over a common
decision horizon. The profiled cost $c$ measures the additional service
work if the object leaves $j$: storage access after DRAM demotion, or
reconstruction after final storage removal. The reuse term multiplies a causal next-use estimate by the corresponding
saved service work, independently
of slack. This preserves useful state for long-slack, multi-round Agents.
When recovery requires a dependent group, its profiled cost is distributed
by bytes rather than repeatedly assigned to every block. Summed Agent
shares express competing interests, not exact joint WCT: simultaneous
delays at a workflow join are not additive.

A one-time grant becomes stale; granting fresh credit at every checkpoint
can retain unused state indefinitely. Inspired by Landlord and
GreedyDual-Size~\cite{young2002landlord,cao1997greedydual}, Master separates
current grant $g$ from already-paid rent $d$:
\begin{equation}
 C_{x,j}=\sum_a[g_{a,x,j}-d_{a,x,j}]^+,
 \qquad H_{x,j}=\rho_j+C_{x,j}/b_x.
 \label{eq:rental-credit}
\end{equation}
The pool pressure clock $\rho_j$ charges bytes $b_x$, rather than elapsed
time. Increasing it by $\delta$ removes $\min(C_{x,j},b_x\delta)$, split
proportionally among positive Agent shares. All committed residents age,
including leased objects; eviction eligibility is separate. Actual reuse or recovery consumption renews only the using Agent,
once per object and request/use ID. Attention observations refresh demand;
checkpoint growth, score updates and migration preserve paid rent.
Expired Agent interest is removed. Figure~\ref{fig:renewable-credit}
illustrates these different events. Recency enters through renewal,
demand through selection probability and reuse history, and workflow
urgency through $u_r$; these are explicit signals, not a claim to subsume
every LRU/LFU variant.

\paragraph{Avoiding a credit update per token.}
The engine advances one query counter per layer and records only selected
block IDs. A block's birth index and the window boundary determine its
opportunities without incrementing every block's denominator. Bounded
query buckets expire old selected counts; dense attention needs no such
counters. Statistics are batched at checkpoint events. Slack changes
update a request scalar and version, leaving its blocks untouched.

Master materializes a block only on use, candidate inspection or a
rotating scan. It first settles rent since that block's last clock,
then reads current request/demand versions and refreshes the grant while
retaining paid amounts. Grants remain fixed between materializations;
the shared pressure clock therefore avoids scanning all residents on
each charge. Refreshing only a stale heap minimum is insufficient when
unseen keys can decrease. Each round refreshes $q$ low-key candidates
and $q$ rotating alternatives, within a fixed work budget, before choosing
victims. This is approximate selection, compared with eager refresh in
Section~\ref{sec:eval-rental}; it is not an exact global minimum. Work
scales with observed selections, examined candidates and their Agent
shares, rather than a full-cache update per generated token.

\subsubsection{Greedy Peer Protection}
\label{sec:sparse-placement}

Peer accelerates recovery without participating in rent. For each supported
model/layout, calibration fixes checkpoint interval $Z$, prefix depth $L$
and $K$ blocks per selected layer. The online policy changes hot-block
identities and request membership, not these three parameters. A Peer
slot reserves at most $LK$ blocks plus a separately budgeted mandatory
state allowance. Short histories use fewer blocks; shared physical
objects are charged once. This is a fixed quota, not a claim of identical
physical occupancy for every request. Engine block size is configured
separately from $Z$; a measurement's 64-token unit is not universal.

Figure~\ref{fig:sparse-recovery} shows the resulting placement. Each sealed
layer/block/TP object has one managed owner, distinct from execution
working KV. Peer, DRAM and storage together cover a committed checkpoint;
moving an object to Peer releases its lower-tier copy after handoff.
Unselected blocks stay below and demand misses fetch their actual owners.

\begin{figure}[t]
 \centering
 \includegraphics[width=\columnwidth]{figures/sparse-recovery.pdf}
 \caption{A checkpoint across three exclusive tiers, with a fixed Peer
 quota $L=2$, $K=2$. Each cell groups a layer/block's TP objects; four
 cells reside in Peer and the others in DRAM/storage. Layer 1 needs
 B0/B1/B2: B0 is local, B1 needs H2D, and B2 needs a storage read plus
 H2D. The dashed tail is uncommitted. Hot identities can change within
 the quota; a Peer victim must migrate down before its space is reusable.}
 \label{fig:sparse-recovery}
 \Description{Three-layer ownership map with two Peer blocks in each of
 the first two layers and actual-owner recovery paths.}
\end{figure}

Only requests with clipped guarded slack $s_r=0$ enter the Peer candidate
queue, including requests whose slack reaches zero while running.
Admitted requests select $K$ blocks per layer by $\widehat p$, with recency
then block ID breaking ties. Completion or tool handoff makes that
request's slot evictable. On admission or quota shrink, Master first
reclaims ended requests, then active requests in decreasing slack order.
Ordinary admission preserves an incumbent at equal slack; a new tied
request waits for Peer while retaining lower-tier protection. Required
shrink breaks active ties by least recent request use. Active requests
therefore have a policy preference, not an absolute eviction lock;
physical read/copy leases are always respected. A shared object is released
only when no retained slot needs it. Different takeover DPs cannot both
count its sole owner as local.

\paragraph{Why fix a partial prefix of layers?}
Let $M$ be committed history and $N$ its acknowledged, uncommitted tail.
Recovery replays $N$; that computation can overlap later-layer reads.
Peer removes only exposed transfer waiting, beyond the reconstruction
already avoided by lower-tier state. Figure~\ref{fig:design-pipeline}
includes storage reads, H2D and replay on one clock. Demand misses make
sparse loading slower than full-layer pipelining in this example; Peer
removes about 15\,ms relative to the same sparse demand. The comparison
illustrates partial residency, not a universally optimal $L$.
Calibration sweeps fixed quotas over declared history/tail and load
profiles, selects a configuration under an HBM budget, then freezes it
before evaluation on other workflows. Heterogeneous requests may prefer
different depths; adaptive-depth and small-instance oracle comparisons
measure the cost of this simplification.

\begin{figure}[t]
 \centering
 \includegraphics[width=\columnwidth]{figures/recovery-pipeline.pdf}
 \caption{Analytical storage/H2D/replay schedules: 32 layers, 32K history,
 64-token tail, and 1,024\,MiB unique KV. (a)--(c) split history equally
 between DRAM/storage; (d) moves $L=7$, $K=336$ 64-token blocks per layer
 to Peer: 147\,MiB Peer, 438.5\,MiB per lower tier. The existing
 MiniCPM4-8B replay calibration is 2.870\,ms/layer; storage/H2D rates
 of 3.125/6.25\,GB/s are assumed. Both sparse cases preload 336 blocks
 and miss on 64 per layer; misses become known after the previous layer
 finishes. These demands are constructed. Cold KV follows the plotted
 endpoint. Placement is complete; I/O stages are independent, workspace
 sufficient and foreground traffic absent. These are mechanism schedules,
 not integrated measurements or an optimal quota search.}
 \label{fig:design-pipeline}
 \Description{Four analytical recovery schedules separating storage
 reads, H2D, computation and exposed wait at equal unique coverage.}
\end{figure}

\subsubsection{Online Three-Tier Management}
\label{sec:placement-plan}

Prefill offers sealed blocks; decode reviews every fixed $Z$ tokens, and
completion/tool handoff offers a final snapshot, including short rounds.
Actual use, Slack changes, capacity updates and transfer outcomes also
invoke Algorithm~\ref{alg:placement}. New blocks receive the creating
Agent's initial grant; unchanged prefix objects receive no new grant.

\paragraph{Admission and pressure.}
Free lower-tier space admits new state, including zero-credit state.
Under pressure, Master refreshes candidate victims and trials enough
byte-proportional rent to exhaust an examined, safely removable bundle.
Incoming objects pay the same trial rent; those reaching zero are deferred
from that tier. DRAM rejection can send state directly to storage.
Committed zero-credit objects remain available until space is needed;
DRAM victims migrate to storage, whose separate ledger values retention
against reconstruction. Last-copy removal must publish a valid older
checkpoint or input/log replay fallback for every affected manifest.
Dependencies and physical leases constrain all victim choices.

Peer admission/shrink follows the request order above. A hot-set change
swaps identities within the same slot. Both can require Peer-to-DRAM and
DRAM-to-storage cascades. If a safe destination cannot be reserved, the
move waits and the source remains occupied. Paid amounts persist per
lower pool across round trips, with no rent during absence. Neither a
move through Peer nor a value-only refresh renews credit.

\par\medskip\noindent\begin{minipage}{\columnwidth}
\raggedright
 \refstepcounter{designalgorithm}\label{alg:placement}
 \textbf{Algorithm \thedesignalgorithm: Online three-tier management}
 \par\smallskip
 \begin{tabularx}{\columnwidth}{@{}r@{\quad}X@{}}
 \toprule
 1 & \textsc{ReconcileTransfers}$(e)$; \textsc{PublishSignals}$(e)$ \\
 2 & \textsc{RecordNewBlocksAndUses}$(e)$ \\
 3 & $Q\gets\textsc{LowKeysAndRotation}(q)$ \\
 4 & \textbf{for} $x\in Q$ \textbf{while} \textsc{WorkLeft}() \textbf{do} \\
 5 & \quad \textsc{SettleRent}$(x)$; \textsc{RefreshGrantAndKey}$(x)$ \\
 6 & $P_D\gets\textsc{LowerTierAdmission}(e,Q)$ \\
 7 & $P_H\gets\textsc{PeerSlots}(s,B,L,K,e)$ \\
 8 & $P\gets\textsc{PrepareCascade}(P_D,P_H)$ \\
 9 & \textbf{if} $P=\varnothing$ \textbf{then} \textsc{Defer}$(e)$; \textbf{return} \\
 10 & \textbf{if not} \textsc{ReserveIfCurrent}$(P)$ \textbf{then} \\
 11 & \quad \textsc{Requeue}$(e)$; \textbf{return} \\
 12 & \textsc{Submit}$(P)$ \quad\textit{completion/failure emits a new event} \\
 \bottomrule
 \end{tabularx}
\end{minipage}\par\medskip

Connector copies, confirms the destination, redirects ownership, and
releases the source after readers drain. Both copies count during the
handoff. Rejected previews change neither ownership nor committed rent;
successful capacity handoffs commit their trial rent and residual grants.
Only complete dependency coverage advances $M$; copy lag can make $N>Z$.
The protected source failure must leave advertised owners reachable;
additional owner loss requires an explicit fallback. Capacity checks
validate feasible execution, not global WCT optimality. Dynamic values,
dependent objects and budgeted refresh also prevent inheriting Landlord's
competitive guarantee. Section~\ref{sec:eval-rental} compares against the
same values with direct revaluation and isolates renewal, aging and
shared-Agent contributions.

\begin{figure}[t]
 \centering
 \includegraphics[width=\columnwidth]{figures/renewable-credit.pdf}
 \caption{Credit for one shared, unit-size DRAM object. Grants 3 and 2
 become remaining shares 1.8 and 1.2 after rent of 2. Updating A's
 grant to 4 preserves its paid amount 1.2, leaving 2.8; an actual new
 use by A renews its share to 4. B stays at 1.2. Lazy materialization
 settles earlier rent before updating the grant. This is a constructed
 accounting example; Peer has no rental clock.}
 \label{fig:renewable-credit}
 \Description{A per-Agent ledger separating aging, value refresh and renewal.}
\end{figure}

